\documentclass[12pt]{article}
\usepackage{amsmath, amssymb}
\usepackage{geometry}
\usepackage{booktabs}
\usepackage{array}
\usepackage{enumitem}
\usepackage{titlesec}
\usepackage[dvipsnames, table]{xcolor}
\usepackage{fancyhdr}
\usepackage{tcolorbox}
\usepackage{microtype}

\tcbuselibrary{skins, breakable}

\geometry{margin=1in, top=1in, bottom=1in}

% ── Colour palette (Charcoal + Cyan + Orange — Material Design) ─────────────
\definecolor{primary}  {RGB}{38,  50,  56}   % dark blue-grey charcoal
\definecolor{accent}   {RGB}{0,  188, 212}   % vivid cyan
\definecolor{lightbg}  {RGB}{224,247,250}    % pale cyan
\definecolor{gold}     {RGB}{255,152,  0}    % vibrant orange
\definecolor{goldbg}   {RGB}{255,243,224}    % light orange cream
\definecolor{softgray} {RGB}{245,246,247}
\definecolor{darktext} {RGB}{33,  33,  33}
\definecolor{green1}   {RGB}{0,  150, 136}   % teal
\definecolor{greenbg}  {RGB}{224,242,241}    % light teal

% ── Page headers / footers ───────────────────────────────────────────────────
\pagestyle{fancy}
\fancyhf{}
\fancyhead[L]{\small\color{primary}\textbf{Operations Research}\;\textcolor{accent}{|}\;Mid Term Solutions}
\fancyhead[R]{\small\color{darktext}BS P-IV Mathematics}
\renewcommand{\headrulewidth}{0pt}
\fancyhead[C]{\color{accent}\rule[-1ex]{\textwidth}{1.2pt}}
\fancyfoot[C]{\small\color{darktext}\thepage}

% ── Section styling ──────────────────────────────────────────────────────────
\titleformat{\section}
  {\large\bfseries\color{primary}}
  {}
  {0em}
  {}
  [\color{accent}\rule{\linewidth}{2pt}]
\titlespacing*{\section}{0pt}{20pt}{10pt}

\titleformat{\subsection}
  {\normalsize\bfseries\color{accent}}
  {}
  {0em}
  {}

% ── tcolorbox styles ─────────────────────────────────────────────────────────
\newtcolorbox{formulabox}{
  enhanced,
  colback=goldbg, colframe=gold, arc=5pt, boxrule=1.5pt,
  left=12pt, right=12pt, top=10pt, bottom=10pt,
  title={\small\bfseries Formula / Key Result},
  colbacktitle=gold!30, coltitle=darktext,
  attach boxed title to top left={xshift=12pt, yshift=-\tcboxedtitleheight/2},
  boxed title style={arc=3pt, colframe=gold, colback=gold!25, boxrule=1pt},
}

\newtcolorbox{problembox}{
  enhanced,
  colback=lightbg, colframe=primary, arc=5pt, boxrule=2pt,
  left=12pt, right=12pt, top=8pt, bottom=8pt,
  title={\small\bfseries Question},
  colbacktitle=primary, coltitle=white,
}

\newtcolorbox{conceptbox}{
  enhanced,
  borderline west={4pt}{0pt}{accent},
  colback=lightbg!55, colframe=white,
  arc=0pt, boxrule=0pt,
  left=12pt, right=8pt, top=6pt, bottom=6pt,
}

\newtcolorbox{observebox}{
  enhanced,
  colback=softgray, colframe=primary!50, arc=5pt, boxrule=1pt,
  left=12pt, right=12pt, top=8pt, bottom=8pt,
  title={\small\bfseries Key Points},
  colbacktitle=primary, coltitle=white,
}

\newtcolorbox{answerbox}{
  enhanced,
  colback=greenbg, colframe=green1, arc=5pt, boxrule=1.5pt,
  left=12pt, right=12pt, top=8pt, bottom=8pt,
  title={\small\bfseries Answer},
  colbacktitle=green1!80, coltitle=white,
  attach boxed title to top left={xshift=12pt, yshift=-\tcboxedtitleheight/2},
  boxed title style={arc=3pt, colframe=green1, colback=green1!70, boxrule=1pt},
}

\newcolumntype{C}[1]{>{\centering\arraybackslash}p{#1}}

\newcommand{\hcell}[1]{\textbf{\color{white}#1}}

\begin{document}

% ════════════════════════════════════════════════════════════════════════════
%  TITLE BANNER
% ════════════════════════════════════════════════════════════════════════════
\thispagestyle{empty}

\begin{tcolorbox}[
  enhanced, arc=0pt, boxrule=0pt,
  colback=primary, colframe=primary,
  left=16pt, right=16pt, top=18pt, bottom=14pt,
  borderline south={4pt}{0pt}{accent},
]
  \centering
  {\huge\bfseries\color{white} Operations Research}\\[4pt]
  {\large\color{accent} Mid-Term Examinations 2026 --- Step-by-Step Solutions}\\[6pt]
  {\normalsize\color{lightbg}
    Shah Abdul Latif University, Khairpur \quad\textbullet\quad Department of Mathematics}
\end{tcolorbox}

\vspace{4pt}

\begin{tcolorbox}[
  enhanced, arc=0pt, boxrule=0pt,
  colback=white, colframe=white,
  borderline north={1.5pt}{0pt}{accent!60},
  borderline south={1.5pt}{0pt}{accent!60},
  left=16pt, right=16pt, top=12pt, bottom=12pt,
]
  \begin{tabular}{p{3.5cm} p{9cm}}
    \textcolor{primary}{\textbf{Class / Session:}}
      & \textbf{BS-IV, Semester: First, Session: 2026} \\[4pt]
    \textcolor{primary}{\textbf{Date / Time:}}
      & \textbf{05-03-2026} \quad|\quad Total Marks: \textbf{30} \quad|\quad Time: \textbf{60 minutes} \\[4pt]
    \textcolor{primary}{\textbf{Instructor:}}
      & \textbf{Sir Darshan Lal} \\[4pt]
    \textcolor{primary}{\textbf{Student:}}
      & \textbf{Nadir Ali Khan} \quad|\quad Roll No: \textbf{52} \\
  \end{tabular}
\end{tcolorbox}

\vspace{14pt}

% ════════════════════════════════════════════════════════════════════════════
\section{Question No.\ 1 --- Graphical Analysis (k-parameter LP)}
% ════════════════════════════════════════════════════════════════════════════

\begin{problembox}
\textbf{Maximize} $Z = x_1 + 2x_2$\\[4pt]
Subject to:
\begin{align*}
  x_1 &\leq 2 \tag{C1}\\
  x_2 &\leq 3 \tag{C2}\\
  kx_1 + x_2 &\leq 2k + 3, \quad k \geq 0 \tag{C3}\\
  x_1, x_2 &\geq 0
\end{align*}
The current solution is $x_1 = 2,\; x_2 = 3$.
Find all values of $k$ for which this solution remains \textbf{optimal}.
\end{problembox}

\subsection*{Step 1 --- Feasibility Check of $(2,3)$}

\begin{conceptbox}
We verify that $(2,3)$ satisfies all constraints for every $k \geq 0$.
\end{conceptbox}

\begin{itemize}[noitemsep, topsep=4pt]
  \item \textbf{C1} ($x_1 \leq 2$): $2 \leq 2$ \quad\checkmark\quad (binding)
  \item \textbf{C2} ($x_2 \leq 3$): $3 \leq 3$ \quad\checkmark\quad (binding)
  \item \textbf{C3} ($kx_1 + x_2 \leq 2k+3$): $k(2)+3 = 2k+3 \leq 2k+3$ \quad\checkmark\quad (binding for all $k \geq 0$)
\end{itemize}

Hence $(2,3)$ is \textbf{always feasible} for all $k \geq 0$.
The objective value at this point is:
\[
  Z = 2 + 2(3) = \mathbf{8}
\]

\subsection*{Step 2 --- Corner Points and Objective Values}

The feasible region is determined by C1 and C2 (with non-negativity), giving the rectangle
$[0,2]\times[0,3]$.  The four corner points and their $Z$-values are:

\begin{center}
\begin{tabular}{C{2.8cm} C{1.6cm} C{1.6cm} C{3.2cm}}
\toprule
\rowcolor{primary}
\hcell{Corner Point} & \hcell{$x_1$} & \hcell{$x_2$} & \hcell{$Z = x_1 + 2x_2$} \\
\midrule
\rowcolor{lightbg!50}
$(0,0)$ & 0 & 0 & 0 \\
$(2,0)$ & 2 & 0 & 2 \\
\rowcolor{lightbg!50}
$(0,3)$ & 0 & 3 & 6 \\
$(2,3)$ & 2 & 3 & $\mathbf{8}$ \quad $\leftarrow$ Maximum \\
\bottomrule
\end{tabular}
\end{center}

\subsection*{Step 3 --- Redundancy of Constraint C3}

Rewrite C3 by rearranging terms:
\[
  kx_1 + x_2 \leq 2k + 3
  \quad\Longleftrightarrow\quad
  k(x_1 - 2) \leq 3 - x_2
\]

\begin{itemize}[noitemsep, topsep=4pt]
  \item From C1: $x_1 \leq 2$, so $x_1 - 2 \leq 0$, hence $k(x_1-2) \leq 0$ for $k \geq 0$.
  \item From C2: $x_2 \leq 3$, so $3 - x_2 \geq 0$.
  \item Therefore $k(x_1-2) \leq 0 \leq 3 - x_2$ \quad for all $k \geq 0$.
\end{itemize}

\begin{conceptbox}
C3 is \textbf{automatically satisfied} at every point in the feasible region defined by C1 and C2.
The feasible region remains the rectangle $[0,2]\times[0,3]$ for \emph{all} $k \geq 0$,
and the corner point $(2,3)$ is never removed from the feasible region.
\end{conceptbox}

\subsection*{Step 4 --- Optimality Conclusion}

The objective function gradient direction is $(1,\,2)$.  Moving the isoprofit line
$Z = x_1 + 2x_2 = c$ outward from the origin, the \textbf{last feasible corner reached}
is always $(2,3)$ with $Z = 8$, regardless of $k$.

The isoprofit line $x_1 + 2x_2 = 8$ touches the feasible region \emph{exactly} at the corner
$(2,3)$ for every $k \geq 0$ because C3 is redundant and never cuts off this corner.

\begin{formulabox}
\textbf{Conclusion:}\\[4pt]
The solution $(x_1,\, x_2) = (2,\, 3)$ with $Z = 8$ is \textbf{optimal for ALL} $k \geq 0$.\\[4pt]
C3 is a redundant constraint for the entire feasible rectangle $[0,2]\times[0,3]$ whenever $k \geq 0$.
\end{formulabox}

\newpage

% ════════════════════════════════════════════════════════════════════════════
\section{Question No.\ 2 --- Metalco Company LP Formulation}
% ════════════════════════════════════════════════════════════════════════════

\begin{problembox}
The Metalco Company must blend a new alloy with the following composition requirements:\\[4pt]
\quad\textbf{40\% tin}, \quad\textbf{35\% zinc}, \quad\textbf{25\% lead}\\[6pt]
Five alloys are available with the following properties:

\vspace{8pt}
\begin{center}
\begin{tabular}{lC{1.8cm}C{1.8cm}C{1.8cm}C{1.8cm}C{1.8cm}}
\toprule
\rowcolor{primary}
\hcell{Property} & \hcell{Alloy 1} & \hcell{Alloy 2} & \hcell{Alloy 3} & \hcell{Alloy 4} & \hcell{Alloy 5} \\
\midrule
\rowcolor{lightbg!50}
Tin (\%)       & 60 & 25 & 45 & 20 & 50 \\
Zinc (\%)      & 10 & 15 & 45 & 50 & 40 \\
\rowcolor{lightbg!50}
Lead (\%)      & 30 & 60 & 10 & 30 & 10 \\
Cost (\$/lb)   & 22 & 20 & 25 & 20 & 27 \\
\bottomrule
\end{tabular}
\end{center}

\vspace{4pt}
Formulate a \textbf{Linear Programming} model to minimize the cost per pound of the new alloy.
\end{problembox}

\subsection*{Decision Variables}

\begin{conceptbox}
Let $x_j$ denote the \textbf{fraction of alloy $j$} ($j = 1, 2, 3, 4, 5$) used in the blend,
so that $0 \leq x_j \leq 1$ for each $j$.
\end{conceptbox}

\subsection*{Objective Function --- Minimize Cost per Pound}

\[
  \text{Min}\quad Z = 22x_1 + 20x_2 + 25x_3 + 20x_4 + 27x_5
\]

\subsection*{Constraints}

\textbf{(i) Blending constraint} (total fractions must sum to 1):
\[
  x_1 + x_2 + x_3 + x_4 + x_5 = 1
\]

\textbf{(ii) Tin requirement} (40\% of the blend):
\[
  60x_1 + 25x_2 + 45x_3 + 20x_4 + 50x_5 = 40
\]

\textbf{(iii) Zinc requirement} (35\% of the blend):
\[
  10x_1 + 15x_2 + 45x_3 + 50x_4 + 40x_5 = 35
\]

\textbf{(iv) Lead requirement} (25\% of the blend):
\[
  30x_1 + 60x_2 + 10x_3 + 30x_4 + 10x_5 = 25
\]

\textbf{(v) Non-negativity:}
\[
  x_j \geq 0, \quad j = 1, 2, 3, 4, 5
\]

\begin{observebox}
\textbf{Note:} The sum constraint $\sum x_j = 1$ is technically implied because
tin\% + zinc\% + lead\% = 100 for each alloy, and $40 + 35 + 25 = 100$.
However, it is retained explicitly for completeness and to ensure the fractions
sum to unity when solving numerically.
\end{observebox}

\vspace{8pt}

\begin{formulabox}
\textbf{Complete LP Model:}\\[6pt]
\[
  \text{Min}\quad Z = 22x_1 + 20x_2 + 25x_3 + 20x_4 + 27x_5
\]
Subject to:
\begin{align*}
  x_1 + x_2 + x_3 + x_4 + x_5 &= 1 \quad\text{(blending)}\\
  60x_1 + 25x_2 + 45x_3 + 20x_4 + 50x_5 &= 40 \quad\text{(tin)}\\
  10x_1 + 15x_2 + 45x_3 + 50x_4 + 40x_5 &= 35 \quad\text{(zinc)}\\
  30x_1 + 60x_2 + 10x_3 + 30x_4 + 10x_5 &= 25 \quad\text{(lead)}\\
  x_j &\geq 0, \quad j = 1,\ldots,5
\end{align*}
\end{formulabox}

\newpage

% ════════════════════════════════════════════════════════════════════════════
\section{Question No.\ 3 --- Standard Form and Initial Simplex Tableau}
% ════════════════════════════════════════════════════════════════════════════

\begin{problembox}
Consider the following LP problem:\\[6pt]
\textbf{Minimize}\quad $z = 2x_1 + 2x_2 - 4x_3$\\[4pt]
Subject to:
\begin{align*}
  2x_1 + 2x_2 + 2x_3 &= 10 \tag{1}\\
  -2x_1 + 6x_2 - x_3 &\leq -10 \tag{2}\\
  x_1 + 3x_2 &\geq 3 \tag{3}\\
  x_1, x_2, x_3 &\geq 0
\end{align*}
\textbf{(a)} Convert to standard form (Big-M method).\\
\textbf{(b)} Write the initial Simplex Tableau-1 and identify the entering and leaving variables.
\end{problembox}

\subsection*{Part (a) --- Standard Form}

We introduce surplus and artificial variables as required by the Big-M method.

\medskip
\noindent\textbf{Constraint (1) --- Equality:}\\
Already an equality. Add artificial variable $A_1 \geq 0$:
\[
  2x_1 + 2x_2 + 2x_3 + A_1 = 10
\]

\medskip
\noindent\textbf{Constraint (2) --- Negative RHS ($\leq$ with $-10$):}\\
Negative RHS: multiply through by $-1$ to reverse the inequality:
\[
  2x_1 - 6x_2 + x_3 \geq 10
\]
Subtract surplus variable $s_1 \geq 0$, add artificial $A_2 \geq 0$:
\[
  2x_1 - 6x_2 + x_3 - s_1 + A_2 = 10
\]

\medskip
\noindent\textbf{Constraint (3) --- $\geq$ type:}\\
Subtract surplus variable $s_2 \geq 0$, add artificial $A_3 \geq 0$:
\[
  x_1 + 3x_2 - s_2 + A_3 = 3
\]

\medskip
\noindent\textbf{Big-M Objective (minimization):}\\
Artificial variables enter the objective with penalty coefficient $M > 0$:
\[
  \text{Min}\quad z = 2x_1 + 2x_2 - 4x_3 + 0{\cdot}s_1 + 0{\cdot}s_2 + M{\cdot}A_1 + M{\cdot}A_2 + M{\cdot}A_3
\]

\begin{formulabox}
\textbf{Complete Standard Form:}\\[6pt]
\[
  \text{Min}\quad z = 2x_1 + 2x_2 - 4x_3 + M A_1 + M A_2 + M A_3
\]
Subject to:
\begin{align*}
  2x_1 + 2x_2 + 2x_3 + A_1 &= 10\\
  2x_1 - 6x_2 + x_3 - s_1 + A_2 &= 10\\
  x_1 + 3x_2 - s_2 + A_3 &= 3\\
  x_1, x_2, x_3, s_1, s_2, A_1, A_2, A_3 &\geq 0
\end{align*}
\end{formulabox}

\newpage

\subsection*{Part (b) --- Initial Simplex Tableau-1}

\begin{conceptbox}
\textbf{Initial Basic Feasible Solution (BFS):}
The artificial variables form the initial basis.\\
$A_1 = 10$, $A_2 = 10$, $A_3 = 3$; all other variables $= 0$.
\end{conceptbox}

\medskip
\noindent\textbf{Computing the $z$-row reduced costs:}

Before row operations, the $z$-row coefficient of each basic variable must be made zero.
We subtract $M \times$ (Row 1) $+\; M \times$ (Row 2) $+\; M \times$ (Row 3) from the
objective row.

The reduced costs (net $z$-row entries) are:

\begin{itemize}[noitemsep, topsep=4pt]
  \item $x_1$:\quad $2 - M(2) - M(2) - M(1) = 2 - 5M$
  \item $x_2$:\quad $2 - M(2) - M(-6) - M(3) = 2 + M$
  \item $x_3$:\quad $-4 - M(2) - M(1) - M(0) = -4 - 3M$
  \item $s_1$:\quad $0 - M(0) - M(-1) - M(0) = M$
  \item $s_2$:\quad $0 - M(0) - M(0) - M(-1) = M$
  \item $A_1, A_2, A_3$:\quad $0$ (basic, by construction)
  \item \textbf{RHS}:\quad $0 - 10M - 10M - 3M = -23M$
\end{itemize}

\medskip
\noindent\textbf{Simplex Tableau-1:}

\begin{center}
\setlength{\tabcolsep}{5pt}
\begin{tabular}{C{1cm} C{1.4cm} C{1.4cm} C{1.4cm} C{1.2cm} C{1.2cm} C{1cm} C{1cm} C{1cm} C{1.4cm}}
\toprule
\rowcolor{primary}
\hcell{BV} & \hcell{$x_1$} & \hcell{$x_2$} & \hcell{$x_3$} & \hcell{$s_1$} & \hcell{$s_2$} & \hcell{$A_1$} & \hcell{$A_2$} & \hcell{$A_3$} & \hcell{RHS} \\
\midrule
\rowcolor{lightbg!50}
$A_1$ & 2 & 2 & 2 & 0 & 0 & 1 & 0 & 0 & 10 \\
$A_2$ & 2 & $-6$ & 1 & $-1$ & 0 & 0 & 1 & 0 & 10 \\
\rowcolor{lightbg!50}
$A_3$ & 1 & 3 & 0 & 0 & $-1$ & 0 & 0 & 1 & 3 \\
\midrule
$z$   & $2{-}5M$ & $2{+}M$ & $-4{-}3M$ & $M$ & $M$ & 0 & 0 & 0 & $-23M$ \\
\bottomrule
\end{tabular}
\end{center}

\vspace{10pt}

\subsection*{Identifying Entering and Leaving Variables}

\begin{conceptbox}
For \textbf{minimization} using Big-M, the entering variable is the one with the \textbf{most negative} reduced cost in the $z$-row (since we want to decrease $z$).
\end{conceptbox}

\medskip
\noindent\textbf{Entering Variable Selection:}

The $z$-row reduced costs are:
\[
  x_1: 2-5M \approx -5M \quad(\text{most negative for large } M),
  \quad x_3: -4-3M \approx -3M,
  \quad \text{others positive}
\]

For large $M$, $2 - 5M < -4 - 3M$ since $-5M < -3M$.
Therefore the most negative coefficient is $2 - 5M$.

\medskip
\noindent\textbf{Entering variable: $x_1$}

\medskip
\noindent\textbf{Minimum Ratio Test (Leaving Variable):}

Only rows with a \emph{positive} coefficient in the $x_1$-column are considered:

\begin{center}
\begin{tabular}{C{2.2cm} C{2.5cm} C{2.5cm} C{3cm}}
\toprule
\rowcolor{primary}
\hcell{Row (BV)} & \hcell{RHS} & \hcell{$x_1$-column} & \hcell{Ratio RHS / coeff} \\
\midrule
\rowcolor{lightbg!50}
$A_1$ & 10 & 2 & $10/2 = 5$ \\
$A_2$ & 10 & 2 & $10/2 = 5$ \\
\rowcolor{lightbg!50}
$A_3$ & 3  & 1 & $3/1 = \mathbf{3}$ \quad $\leftarrow$ minimum \\
\bottomrule
\end{tabular}
\end{center}

\textbf{Leaving variable: $A_3$} (row 3, minimum ratio $= 3$). The pivot element is $\mathbf{1}$ (row 3, column $x_1$).

\begin{answerbox}
\textbf{Entering variable:} $x_1$ \quad (column) \;---\; most negative reduced cost $= 2 - 5M$\\[4pt]
\textbf{Leaving variable:} $A_3$ \quad (row 3) \;---\; minimum ratio $= 3/1 = 3$\\[4pt]
\textbf{Pivot element:} $1$ \quad (row $A_3$, column $x_1$)
\end{answerbox}

\newpage

% ════════════════════════════════════════════════════════════════════════════
\section{Summary of Key Results}
% ════════════════════════════════════════════════════════════════════════════

\begin{formulabox}
\textbf{Q1 --- Graphical Analysis (k-parameter LP):}\\[4pt]
The solution $(x_1, x_2) = (2, 3)$ with $Z = 8$ is \textbf{optimal for ALL} $k \geq 0$.\\
Constraint C3 is redundant over the feasible rectangle $[0,2]\times[0,3]$ for every $k \geq 0$,
since $k(x_1 - 2) \leq 0 \leq 3 - x_2$ always holds.\\[10pt]
%
\textbf{Q2 --- Metalco Company LP:}\\[4pt]
\[
  \text{Min}\; Z = 22x_1 + 20x_2 + 25x_3 + 20x_4 + 27x_5
\]
Subject to four equality constraints (blending, tin, zinc, lead) and $x_j \geq 0$.\\[10pt]
%
\textbf{Q3 --- Standard Form \& Simplex Tableau-1:}\\[4pt]
\begin{tabular}{lp{10cm}}
Standard form: & 3 equality constraints with $A_1, A_2, A_3$ (artificials) and $s_1, s_2$ (surplus). \\[4pt]
Initial BFS: & $A_1 = 10,\; A_2 = 10,\; A_3 = 3$; all others $= 0$. \\[4pt]
Pivot: & Entering $x_1$ (reduced cost $2-5M$), Leaving $A_3$ (ratio $= 3$), pivot element $= 1$.
\end{tabular}
\end{formulabox}

\end{document}
